% Generated by scripts/build_mixed.py
\documentclass[landscape,8pt,a4paper,twocolumn]{extarticle}
\usepackage{fontspec}
\usepackage{kotex}
\setmainfont{Malgun Gothic}
\setmonofont{Consolas}
\setmonohangulfont{Malgun Gothic}
\usepackage{ifthen}
\usepackage{teamnote}
\teamnote{ICPC}{C++ / Python Teamnote}{}
\geometry{left=8mm,right=8mm,top=12mm,bottom=8mm}
\setminted{fontsize=\scriptsize,baselinestretch=0.95}
\begin{document}
\tableofcontents
\subsection*{C++17 common definitions}
\inputminted{cpp}{mixed/cpp/common.hpp}
\section{정렬 / 이분 탐색 관련}
\subsection{정렬 비교 함수 구현 [PY]}
\inputminted{python}{python/Sorting/Compare.py}
\subsection{좌표 압축 [PY]}
\inputminted{python}{python/Sorting/CoordCompress.py}
\subsection{이분 탐색과 선택 알고리즘 [PY]}
\inputminted{python}{python/Sorting/STL.py}
\section{그래프/트리 (Easy)}
\subsection{인접 리스트, DFS, BFS [PY]}
\inputminted{python}{python/Graph1/DFSBFS.py}
\subsection{위상 정렬 [PY]}
\inputminted{python}{python/Graph1/TopSort.py}
\subsection{최단 거리 - Floyd Warshall [PY/C++]}
\inputminted{python}{python/Graph1/Floyd.py}
\inputminted{cpp}{mixed/cpp/Graph1/Floyd.hpp}
\subsection{최단 거리 - Dijkstra [PY/C++]}
\inputminted{python}{python/Graph1/Dijkstra.py}
\inputminted{cpp}{mixed/cpp/Graph1/Dijkstra.hpp}
\subsection{최단 거리 - Bellman Ford [PY/C++]}
\inputminted{python}{python/Graph1/Bellman.py}
\inputminted{cpp}{mixed/cpp/Graph1/Bellman.hpp}
\subsection{유니온 파인드 + 최소 신장 트리(Kruskal) [PY/C++]}
\inputminted{python}{python/Graph1/Kruskal.py}
\inputminted{cpp}{mixed/cpp/Graph1/Kruskal.hpp}
\section{자료구조}
\subsection{세그먼트 트리 [PY/C++]}
\inputminted{python}{python/DataStructure/SegmentTree.py}
\inputminted{cpp}{mixed/cpp/DataStructure/SegmentTree.hpp}
\subsection{세그먼트 트리 + 레이지 프로퍼게이션 [C++]}
\inputminted{cpp}{mixed/cpp/DataStructure/SegmentTreeLazy.hpp}
\subsection{Convex Hull Trick [C++]}
\inputminted{cpp}{mixed/cpp/DataStructure/ConvexHullTrick.hpp}
\subsection{퍼시스턴트 세그먼트 트리 [C++]}
\inputminted{cpp}{mixed/cpp/DataStructure/PersistentSegmentTree.hpp}
\section{그래프/트리 (Hard)}
\subsection{SCC - Kosaraju [PY/C++]}
\inputminted{python}{python/Graph2/SCC.py}
\inputminted{cpp}{mixed/cpp/Graph2/SCC.hpp}
\subsection{BCC - Tarjan [C++]}
\inputminted{cpp}{mixed/cpp/Graph2/BCC.hpp}
\subsection{최대 유량 - Dinic [C++]}
\inputminted{cpp}{mixed/cpp/Graph2/Dinic.hpp}
\subsection{MCMF [C++]}
\inputminted{cpp}{mixed/cpp/Graph2/MCMF.hpp}
\subsection{이분 매칭 - Hopcroft Karp [C++]}
\inputminted{cpp}{mixed/cpp/Graph2/Hopcroft.hpp}
\subsection{최소 공통 조상(LCA) [PY/C++]}
\inputminted{python}{python/Graph2/LCA.py}
\inputminted{cpp}{mixed/cpp/Graph2/LCA.hpp}
\subsection{Heavy Light Decomposition [C++]}
\inputminted{cpp}{mixed/cpp/Graph2/HLD.hpp}
\section{수학}
\subsection{나눗셈, 최대공약수, 최소공배수 [PY]}
\inputminted{python}{python/Math/Division.py}
\subsection{빠른 거듭제곱 [PY]}
\inputminted{python}{python/Math/PowMod.py}
\subsection{소수 판별, 소인수분해 [PY]}
\inputminted{python}{python/Math/Primes.py}
\subsection{에라토스테네스의 체, 소인수분해 [PY]}
\inputminted{python}{python/Math/Sieve.py}
\subsection{선형 시간 체, 곱셈적 함수 전처리 [C++]}
\inputminted{cpp}{mixed/cpp/Math/LinearSieve.hpp}
\subsection{확장 유클리드 알고리즘 [PY]}
\inputminted{python}{python/Math/ExtendedEuclidean.py}
\subsection{중국인의 나머지 정리 [PY/C++]}
\inputminted{python}{python/Math/ChineseRemainderTheorem.py}
\inputminted{cpp}{mixed/cpp/Math/ChineseRemainderTheorem.hpp}
\subsection{이항 계수를 소수로 나눈 나머지 [PY]}
\inputminted{python}{python/Math/LucasTheorem.py}
\subsection{빠른 소수 판별, 소인수분해 - Miller Rabin, Pollard Rho [C++]}
\inputminted{cpp}{mixed/cpp/Math/MillerRabin-PollardRho.hpp}
\subsection{가우스 소거법 - RREF, 랭크, 행렬식, 역행렬 [PY/C++]}
\inputminted{python}{python/Math/Matrix.py}
\inputminted{cpp}{mixed/cpp/Math/Matrix.hpp}
\subsection{다항식 곱셈(FFT) [C++]}
\inputminted{cpp}{mixed/cpp/Math/Convolution.hpp}
\section{문자열}
\subsection{문자열 해싱 [C++]}
\inputminted{cpp}{mixed/cpp/String/Hashing.hpp}
\subsection{문자열 매칭 - KMP [PY/C++]}
\inputminted{python}{python/String/KMP.py}
\inputminted{cpp}{mixed/cpp/String/KMP.hpp}
\subsection{가장 긴 팰린드롬 부분 문자열 - Manacher [PY/C++]}
\inputminted{python}{python/String/Manacher.py}
\inputminted{cpp}{mixed/cpp/String/Manacher.hpp}
\subsection{문자열 매칭 - Z [PY/C++]}
\inputminted{python}{python/String/Z.py}
\inputminted{cpp}{mixed/cpp/String/Z.hpp}
\subsection{접미사 배열 [C++]}
\inputminted{cpp}{mixed/cpp/String/SuffixArray.hpp}
\section{계산 기하}
\subsection{2차원 계산 기하 템플릿 + CCW [PY/C++]}
\inputminted{python}{python/Geometry/Template.py}
\inputminted{cpp}{mixed/cpp/Geometry/Template.hpp}
\subsection{360도 각도 정렬 [PY/C++]}
\inputminted{python}{python/Geometry/PolarSort.py}
\inputminted{cpp}{mixed/cpp/Geometry/PolarSort.hpp}
\subsection{다각형 넓이 [PY]}
\inputminted{python}{python/Geometry/PolygonArea.py}
\subsection{선분 교차 판정 [PY/C++]}
\inputminted{python}{python/Geometry/SegmentIntersection.py}
\inputminted{cpp}{mixed/cpp/Geometry/SegmentIntersection.hpp}
\subsection{다각형 내부 판별 [PY/C++]}
\inputminted{python}{python/Geometry/PointInPolygon.py}
\inputminted{cpp}{mixed/cpp/Geometry/PointInPolygon.hpp}
\subsection{볼록 껍질 - Monotone Chain [C++]}
\inputminted{cpp}{mixed/cpp/Geometry/ConvexHull.hpp}
\subsection{가장 먼 두 점 - Rotating Calipers [C++]}
\inputminted{cpp}{mixed/cpp/Geometry/Calipers.hpp}
\subsection{볼록 다각형 내부 판별 [C++]}
\inputminted{cpp}{mixed/cpp/Geometry/PointInConvexPolygon.hpp}
\section{기타}
\subsection{가장 긴 증가하는 부분 수열(LIS) [PY/C++]}
\inputminted{python}{python/Misc/LIS.py}
\inputminted{cpp}{mixed/cpp/Misc/LIS.hpp}
\subsection{이분 탐색 [PY]}
\inputminted{python}{python/Misc/BinarySearch.py}
\subsection{삼분 탐색 [PY]}
\inputminted{python}{python/Misc/TernarySearch.py}
\subsection{Python 랜덤, 순서 통계, 비트마스킹 [PY]}
\inputminted{python}{python/Misc/Cpp-Grammer.py}
\subsection{빠른 입력(Fast Input from STDIN) [PY/C++]}
\inputminted{python}{python/Misc/FastIO.py}
\inputminted{cpp}{mixed/cpp/Misc/FastIO.hpp}
\subsection{구간별 약수 최대 개수, 최대 소수}
\inputminted[fontsize=\scriptsize]{text}{docs/NUMBER_TABLES.txt}
\subsection{카탈란 수, 심슨 적분, 그런디 정리, 픽의 정리, 페르마 포인트, 오일러 정리}
\inputminted{python}{python/Misc/Reference.py}
\noindent 카탈란 수: C(n) = comb(2*n,n) // (n+1). 처음 값은 1, 1, 2, 5, 14, 42, 132, 429, 1430, 4862, 16796, 58786, 208012, 742900. 길이 2n의 올바른 괄호 문자열, 잎 n+1개의 순서 있는 full binary tree, 볼록 (n+2)각형의 삼각 분할을 센다.\par\smallskip
\noindent 심슨 공식: n은 양의 짝수, h=(b-a)/n, x\_i=a+i*h. 적분 근삿값은 h/3 * (f(x\_0)+f(x\_n)+4*sum(f(x\_i), i 홀수)+2*sum(f(x\_i), 0<i<n 짝수)). f의 4계 도함수가 연속이고 절댓값이 M 이하이면 절대 오차는 M*(b-a)*h\textasciicircum{}4/180 이하.\par\smallskip
\noindent Nim: 일반 규칙에서는 모든 더미 크기의 XOR가 0이면 후공 승리, 아니면 선공 승리. Grundy 값은 가능한 다음 상태의 값들의 mex(포함되지 않은 최소 비음수 정수). 독립적인 게임의 합은 각 Grundy 값의 XOR. 유한 비순환 impartial normal-play 게임에 적용한다.\par\smallskip
\noindent Subtraction game: 한 번에 한 더미에서 1..k개를 제거할 때 각 더미 크기를 k+1로 나눈 나머지가 Grundy 값이다. 이 값들의 XOR로 판단한다.\par\smallskip
\noindent Index-k Nim (Moore's Nim): 한 번에 1..k개 더미를 고르고 각 더미에서 양의 개수를 제거한다. 각 이진 자릿수에서 1의 개수를 k+1로 나눈 나머지가 모두 0인 상태가 패배 상태다. 이 나머지 벡터를 일반적인 Grundy 값으로 취급하지 않는다.\par\smallskip
\noindent Misere Nim: 마지막 돌을 가져가는 쪽이 패배한다. 0인 더미를 제외하고 모두 크기 1이면 양의 더미 수가 홀수일 때 후공 승리. 크기 2 이상 더미가 있으면 XOR가 0일 때 후공 승리. 빈 시작 상태는 게임 규칙의 종료 정의를 확인한다.\par\smallskip
\noindent Pick 정리: 단순 격자 다각형에서 넓이 A = 내부 격자점 수 I + 경계 격자점 수 B/2 - 1. 경계 격자점 수는 각 변의 gcd(abs(dx), abs(dy)) 합이다. floor\_sum(n,m,a,b)는 i=0..n-1에 대해 floor((a*i+b)/m)를 합한다. 아래 Reference.py에 반복형 구현을 포함했다.\par\smallskip
\noindent 페르마 포인트: 삼각형 꼭짓점까지 거리 합을 최소화한다. 120도 이상인 각이 있으면 그 꼭짓점이다. 모든 각이 120도 미만이면 각 변 바깥쪽 정삼각형의 새 꼭짓점과 반대 꼭짓점을 잇는 선분들이 만나는 점이다. 최소 거리 합은 sqrt((a*a+b*b+c*c+4*sqrt(3)*A)/2), A는 삼각형 넓이.\par\smallskip
\noindent 오일러 정리: gcd(a,m)=1이면 a\textasciicircum{}phi(m) = 1 (mod m). 이 경우 지수를 phi(m)로 나눠 줄일 수 있다. 서로소가 아닐 때 무조건 지수를 줄이면 틀린다. 충분한 지수 e >= ceil(log2(m)), m>=2에서는 a\textasciicircum{}e = a\textasciicircum{}(e \% phi(m) + phi(m)) (mod m)를 사용할 수 있다. 작은 지수는 원래대로 계산한다.\par\smallskip
\noindent 소수 p에 대해 g\textasciicircum{}0+...+g\textasciicircum{}(p-2)는 g=1 (mod p)이면 -1 (mod p), g가 0도 1도 아니면 0 (mod p)이다. g=0이면 합은 1이므로 별도 처리한다.\par\smallskip
\subsection{경우의 수 - 포함 배제, 스털링 수, 벨 수}
\noindent 구별되는 n개 공을 구별되는 k개 상자에 모두 비지 않게 넣는 경우(전사 함수): sum((-1)\textasciicircum{}(k-i)*comb(k,i)*i\textasciicircum{}n, i=0..k) = k!*S2(n,k). 원문 첫 항목의 '공 구별 X'를 '공 구별 O'로 정정했다. 구별되지 않는 공, 구별되는 양의 상자라면 comb(n-1,k-1)이다.\par\smallskip
\noindent 구별되는 공, 구별되지 않는 k개 비지 않은 상자: 제2종 스털링 수 S2(n,k)=k*S2(n-1,k)+S2(n-1,k-1). S2(0,0)=1, 유효 범위 밖은 0. 포함 배제식은 위 전사 함수 수를 k!로 나눈 값이다. 모듈러 나눗셈에는 역원의 존재 조건을 확인한다.\par\smallskip
\noindent 벨 수: B(n)=sum(S2(n,k), k=0..n). 최대 k개 블록이면 B(n,k)=sum(S2(n,j), j=0..k). B(n)=sum(comb(n-1,i)*B(i), i=0..n-1), B(0)=1. 최대 k개 식은 sum(i\textasciicircum{}n/i! * sum((-1)\textasciicircum{}j/j!, j=0..k-i), i=0..k)로도 정리된다. 이 식의 나눗셈은 유리수 또는 유효한 모듈러 역원으로 계산한다.\par\smallskip
\noindent 완전 순열: D(0)=1, D(1)=0, D(n)=(n-1)*(D(n-1)+D(n-2)).\par\smallskip
\noindent 부호 없는 제1종 스털링 수: c(n,k)=(n-1)*c(n-1,k)+c(n-1,k-1). 부호 있는 제1종 스털링 수: s(n,k)=s(n-1,k-1)-(n-1)*s(n-1,k). 둘 다 (0,0)=1. 원문의 signed 점화식 부호를 정정했다.\par\smallskip
\noindent 정수 n을 정확히 k개의 양의 정수로 분할: p(n,k)=p(n-1,k-1)+p(n-k,k), p(0,0)=1. 전체 분할 수: p(0)=1, p(n)=sum((-1)\textasciicircum{}(k-1)*(p(n-k*(3*k-1)//2)+p(n-k*(3*k+1)//2)), k>=1; 음수 인자의 p는 0. 원문의 부호를 정정했다.\par\smallskip
\noindent 카탈란 등가식: C(n)=comb(2*n,n)-comb(2*n,n+1)=(2*n)!/(n!*(n+1)!). C(0)=1, C(n)=sum(C(i)*C(n-1-i), i=0..n-1) for n>=1.\par\smallskip
\subsection{삼각형의 오심 - 외심, 내심, 무게중심, 수심, 방심}
\noindent 비퇴화 삼각형에서 변 a,b,c는 각각 꼭짓점 A,B,C의 맞은편 길이. 반둘레 p=(a+b+c)/2. 넓이 K=sqrt(p*(p-a)*(p-b)*(p-c)). 외접 반지름 R=a*b*c/(4*K), 내접 반지름 r=K/p.\par\smallskip
\noindent 중선 m\_a=sqrt(2*b*b+2*c*c-a*a)/2. 각 이등분선 길이 t\_a=sqrt(b*c*(1-(a/(b+c))\textasciicircum{}2)). 사인 법칙 sin(A)/a=1/(2*R). 코사인 법칙 a*a=b*b+c*c-2*b*c*cos(A). 탄젠트 법칙 (a+b)/(a-b)=tan((A+B)/2)/tan((A-B)/2), 분모가 0이면 별도 처리.\par\smallskip
\noindent 무게 (alpha,beta,gamma)에 대응하는 점은 (alpha*A+beta*B+gamma*C)/(alpha+beta+gamma). X=b*b+c*c-a*a, Y=a*a+c*c-b*b, Z=a*a+b*b-c*c라 두면 외심의 무게는 (a*a*X,b*b*Y,c*c*Z), 내심은 (a,b,c), 무게중심은 (1,1,1), 수심은 (Y*Z,Z*X,X*Y), A에 대응하는 방심은 (-a,b,c). 직각 삼각형의 수심 공식은 0/0이 될 수 있으므로 직각 꼭짓점을 직접 반환한다.\par\smallskip
\subsection{미적분, 뉴턴 랩슨법}
\noindent 실수 정의역에서 (arcsin x)'=1/sqrt(1-x*x), (arccos x)'=-1/sqrt(1-x*x), (arctan x)'=1/(1+x*x), (tan x)'=1+tan(x)\textasciicircum{}2.\par\smallskip
\noindent 적분 tan(a*x) dx = -ln(abs(cos(a*x)))/a + C. 적분 x*sin(a*x) dx = (sin(a*x)-a*x*cos(a*x))/(a*a)+C. 두 식 모두 a != 0.\par\smallskip
\noindent Newton: x\_next=x-f(x)/f'(x). 초기 추정치와 도함수에 따라 발산·순환할 수 있다. 잔차, 도함수 0, 반복 상한을 확인한다. 이분 탐색처럼 수렴이 보장되는 방법과 혼용할 수 있다.\par\smallskip
\noindent Green 정리: 경계 C에서 integral(L dx + M dy) = 영역 D에서 double\_integral(dM/dx - dL/dy) dx dy. C는 양의 방향의 조각별 매끄러운 단순 폐곡선, L,M은 D를 포함하는 열린 영역에서 연속 편도함수를 가져야 한다.\par\smallskip
\subsection{문제 풀이 체크리스트}
\noindent 비슷한 문제를 풀어본 적이 있는가? 단순한 완전 탐색에서 시작할 수 있는가? 예제를 직접 풀고 그 과정을 수식화할 수 있는가? 문제를 단순화하거나 그림으로 표현할 수 있는가?\par\smallskip
\noindent 문제를 작은 부분으로 분해할 수 있는가? 뒤에서부터 풀 수 있는가? 순서를 강제할 수 있는가? 답을 특정한 정규형으로 제한할 수 있는가?\par\smallskip
\noindent 구간을 통째로 선택하는 문제는 플로우와 구간 자료구조를 고려한다. 원문 힌트 (i,i+1,k,0), (s,e,1,w), (N,T,k,0)는 (시작,끝,용량,비용) 간선 예시다. 모델과 비용 부호는 문제에 맞춰 정한다.\par\smallskip
\noindent a=b를 만족시켜야 한다면 a만 변경, b만 변경, 둘 다 변경, 역방향 변환을 각각 생각한다. 당연하다고 가정한 내용을 다시 점검한다.\par\smallskip
\noindent 확률 문제는 DP와 이분 탐색을 고려한다(원문 예: NYPC 2019 Finals C). 최대·최소 문제는 이분 탐색, 그리디의 교환 논증·접두사 고정, 순서 고정 DP를 검토한다.\par\smallskip
\noindent Python 제출 전: 입력량과 메모리, 정수 연산 비용, 0-based/[l,r) 경계, 빈 입력·중복·연결되지 않은 그래프, 음수 가중치, 재귀 사용 여부, float 오차, 해시 충돌을 확인한다. 로컬 테스트 통과와 온라인 저지 통과를 구분한다. \par\smallskip
\end{document}
